\chapter{Control Flow Instructions}
\label{ch:control-flow}

\section{Overview}

Control-flow instructions change the program counter or compute address-register values used by later control-flow
instructions. The SIRC-1 provides these control-flow operations:

\begin{itemize}
	\item \mnemonic{BRAN} -- Branch by adding a displacement to an address register, usually \reg{p}
	\item \mnemonic{BRSR} -- Branch to subroutine, saving a return address in \reg{l}
	\item \mnemonic{LDEA} -- Load Effective Address into an address register pair
	\item \mnemonic{LJMP} -- Long Jump assembler convenience for \texttt{LDEA p, ...}
	\item \mnemonic{LJSR} -- Long Jump to Subroutine, saving a return address in \reg{l}
\end{itemize}

All encoded control-flow instructions can be conditional. If the condition is false, the instruction is skipped as
described in Chapter~\ref{ch:reading-instructions}.

\section{Legal Forms}

\begin{table}[H]
	\centering
	\scriptsize
	\begin{tabular}{llll}
		\toprule
		\textbf{Instruction} & \textbf{Syntax} & \textbf{Encoding} & \textbf{Notes} \\
		\midrule
		BRAN & \texttt{BRAN \#disp} & 0x1A & PC-relative immediate displacement \\
		BRAN & \texttt{BRAN label} & 0x1A & PC-relative linker-computed displacement \\
		BRAN & \texttt{BRAN (\#disp, addr)} & 0x1A & Address-register-relative immediate displacement \\
		BRAN & \texttt{BRAN (rS, addr)} & 0x1B & Address-register-relative register displacement \\
		BRSR & \texttt{BRSR \#disp} & 0x1E & PC-relative call \\
		BRSR & \texttt{BRSR label} & 0x1E & PC-relative linker-computed call \\
		BRSR & \texttt{BRSR (\#disp, addr)} & 0x1E & Address-register-relative call \\
		BRSR & \texttt{BRSR (rS, addr)} & 0x1F & Address-register-relative register call \\
		LDEA & \texttt{LDEA dest, (\#offset, src)} & 0x18 & Immediate-offset effective address \\
		LDEA & \texttt{LDEA dest, (rS, src)} & 0x19 & Register-offset effective address \\
		LJMP & \texttt{LJMP src} & \texttt{LDEA} meta & Assembles to \texttt{LDEA p, (\#0, src)} \\
		LJMP & \texttt{LJMP src, \#offset} & \texttt{LDEA} meta & Assembles to \texttt{LDEA p, (\#offset, src)} \\
		LJMP & \texttt{LJMP src, rS} & \texttt{LDEA} meta & Assembles to \texttt{LDEA p, (rS, src)} \\
		LJSR & \texttt{LJSR src} & 0x1C & Absolute call through address register \\
		LJSR & \texttt{LJSR src, \#offset} & 0x1C & Immediate-offset absolute call \\
		LJSR & \texttt{LJSR src, rS} & 0x1D & Register-offset absolute call \\
		\bottomrule
	\end{tabular}
	\caption{Legal Control-Flow Instruction Forms}
	\label{tab:control-flow-legal-forms}
\end{table}

\noindent\textbf{Address registers:} \texttt{addr}, \texttt{src}, and \texttt{dest} are address register pairs:
\reg{l}, \reg{a}, \reg{s}, or \reg{p}. The \texttt{BRAN \#disp}, \texttt{BRAN label}, \texttt{BRSR \#disp}, and
\texttt{BRSR label} shorthands use \reg{p} as the base address register.

\noindent\textbf{Shift suffixes:} Control-flow displacements are not shifted. Some instruction encodings can make it look
like a shift suffix on register-displacement branch forms is valid, but the encoded shift is not applied to the displacement
and should not be used.

\section{Common Control-Flow Semantics}

\begin{table}[H]
	\centering
	\begin{tabular}{lp{9cm}}
		\toprule
		\textbf{Topic} & \textbf{Definition} \\
		\midrule
		Condition false & The instruction is skipped and has no side effects. The program counter advances to the next instruction normally. \\
		Address calculation & The selected address register high word supplies the segment. The target low word is the selected address register low word plus the immediate or register displacement. \\
		Branch shorthand & \texttt{BRAN \#disp} and \texttt{BRSR \#disp} are PC-relative forms equivalent to using \reg{p} as the base address register. \\
		Subroutine link & \mnemonic{BRSR} and \mnemonic{LJSR} save the normal next-instruction address in \reg{l} before transferring control. \\
		Status flags & Control-flow instructions preserve all status flags and do not support status update override syntax. \\
		Timing & Control-flow instructions complete in 6 cycles, plus any instruction-fetch wait states. They do not perform a data-memory bus access. \\
		Exceptions & Control-flow address calculation can raise segment-overflow faults when trap-on-address-overflow is enabled. \\
		Privilege & In protected mode, control-flow address-register-pair write-back preserves the destination high word and updates only the low word. Direct writes to high address registers still raise a privilege-violation fault. \\
		Fault side effects & Program-counter, destination address-register, and link-register updates occur after condition evaluation. If a fault aborts execution before those updates, the architectural side effects do not complete. \\
		\bottomrule
	\end{tabular}
	\caption{Common Control-Flow Instruction Semantics}
	\label{tab:control-flow-common-semantics}
\end{table}

\section{Branch Instructions}

\begin{instructionbox}{BRAN -- Branch}
	\textbf{Opcodes:} 0x1A (Immediate Displacement), 0x1B (Register Displacement)

	\textbf{Syntax:}
	\begin{lstlisting}
BRAN #disp                    ; p = p + disp
BRAN label                    ; p = p + label offset
BRAN (#disp, addr)            ; p = addr + disp
BRAN (rS, addr)               ; p = addr + rS
BRAN|cond #disp               ; Conditional PC-relative branch
\end{lstlisting}

	\textbf{Operation:}
	\begin{verbatim}
if condition:
    p.high = addr.high
    p.low  = addr.low + displacement
\end{verbatim}

	\textbf{Description:}

	Branches to a new address by adding a displacement to the selected address register. The shorthand form is
	PC-relative and uses \reg{p} as the base. The full forms can branch relative to any address register pair.

	\textbf{Write-back:} If the condition is true, writes the computed address to \reg{p}.

	\textbf{Flags:} Preserved.

	\textbf{Privilege:} See Table~\ref{tab:control-flow-common-semantics}.

	\textbf{Examples:}
	\begin{lstlisting}
; Unconditional branch forward
BRAN #20                      ; Jump ahead 20 words

; Backward branch (loop)
loop:
    ; ... loop body ...
    BRAN|!= @loop             ; Jump back if not zero

; Branch relative to a table base in a
BRAN (r2, a)                  ; p = a + r2
\end{lstlisting}

\end{instructionbox}

\begin{instructionbox}{BRSR -- Branch to Subroutine}
	\textbf{Opcodes:} 0x1E (Immediate Displacement), 0x1F (Register Displacement)

	\textbf{Syntax:}
	\begin{lstlisting}
BRSR #disp                    ; l = next PC; p = p + disp
BRSR label                    ; l = next PC; p = p + label offset
BRSR (#disp, addr)            ; l = next PC; p = addr + disp
BRSR (rS, addr)               ; l = next PC; p = addr + rS
BRSR|cond #disp               ; Conditional PC-relative call
\end{lstlisting}

	\textbf{Operation:}
	\begin{verbatim}
if condition:
    l = address_of_next_instruction
    p.high = addr.high
    p.low  = addr.low + displacement
\end{verbatim}

	\textbf{Description:}

	Calls a nearby or computed subroutine by saving the return address in the link register and branching to the computed
	target. Return with \mnemonic{RETS}, which assembles to \texttt{LDEA p, (\#0, l)}.

	\textbf{Write-back:} If the condition is true, writes the return address to \reg{l} and the computed target address to
	\reg{p}.

	\textbf{Flags:} Preserved.

	\textbf{Privilege:} See Table~\ref{tab:control-flow-common-semantics}.

	\textbf{Examples:}
	\begin{lstlisting}
; PC-relative call
BRSR function

; Call through a table relative to a
LOAD r1, #8
BRSR (r1, a)

; Return from the subroutine
RETS
\end{lstlisting}

\end{instructionbox}

\section{Load Effective Address}

\begin{instructionbox}{LDEA -- Load Effective Address}
	\textbf{Opcodes:} 0x18 (Immediate Offset), 0x19 (Register Offset)

	\textbf{Syntax:}
	\begin{lstlisting}
LDEA dest, (#offset, src)     ; dest = src + offset
LDEA dest, (rS, src)          ; dest = src + rS
LDEA|cond dest, (#offset, src)
\end{lstlisting}

	\textbf{Operation:}
	\begin{verbatim}
if condition:
    dest.high = src.high
    dest.low  = src.low + offset
\end{verbatim}

	\textbf{Description:}

	Computes an effective address and loads it into the destination address register pair. This is useful for pointer
	arithmetic and for control flow when the destination is \reg{p}.

	\textbf{Write-back:} If the condition is true, writes the computed address to \texttt{dest}.

	\textbf{Flags:} Preserved.

	\textbf{Privilege:} See Table~\ref{tab:control-flow-common-semantics}.

	\textbf{Examples:}
	\begin{lstlisting}
; Pointer arithmetic
LDEA a, (#16, a)              ; a = a + 16

; Computed jump
LDEA p, (r1, p)               ; p = p + r1

; Return from subroutine
LDEA p, (#0, l)               ; p = l
\end{lstlisting}

\end{instructionbox}

\section{Long Jump Instructions}

\begin{instructionbox}{LJMP -- Long Jump}
	\textbf{Assembles to:} \mnemonic{LDEA} with \reg{p} as the destination

	\textbf{Syntax:}
	\begin{lstlisting}
LJMP src                      ; p = src
LJMP src, #offset             ; p = src + offset
LJMP src, rS                  ; p = src + rS
LJMP|cond src                 ; Conditional long jump
\end{lstlisting}

	\textbf{Description:}

	Performs an absolute jump through an address register. \mnemonic{LJMP} is an assembler convenience, not a separate CPU
	opcode; it assembles to the corresponding \mnemonic{LDEA p, ...} form.

	\textbf{Write-back:} If the condition is true, writes the computed address to \reg{p}.

	\textbf{Flags:} Preserved.

	\textbf{Privilege:} Same as \mnemonic{LDEA p, ...}. See Table~\ref{tab:control-flow-common-semantics}.

	\textbf{Examples:}
	\begin{lstlisting}
; Simple absolute jump
LOAD ah, #0x0040
LOAD al, #0x0000              ; a = 0x00400000
LJMP a                        ; Jump to 0x00400000

; Jump with immediate offset
LJMP a, #16                   ; p = a + 16

; Jump with register offset
LOAD r1, #8
LJMP a, r1                    ; p = a + r1
\end{lstlisting}

\end{instructionbox}

\begin{instructionbox}{LJSR -- Long Jump to Subroutine}
	\textbf{Opcodes:} 0x1C (Immediate Offset), 0x1D (Register Offset)

	\textbf{Syntax:}
	\begin{lstlisting}
LJSR src                      ; l = next PC; p = src
LJSR src, #offset             ; l = next PC; p = src + offset
LJSR src, rS                  ; l = next PC; p = src + rS
LJSR|cond src                 ; Conditional long call
\end{lstlisting}

	\textbf{Operation:}
	\begin{verbatim}
if condition:
    l = address_of_next_instruction
    p.high = src.high
    p.low  = src.low + offset
\end{verbatim}

	\textbf{Description:}

	Calls a subroutine through an address register. The return address is saved in \reg{l}; return with \mnemonic{RETS}.

	\textbf{Write-back:} If the condition is true, writes the return address to \reg{l} and the computed target address to
	\reg{p}.

	\textbf{Flags:} Preserved.

	\textbf{Privilege:} See Table~\ref{tab:control-flow-common-semantics}.

	\textbf{Examples:}
	\begin{lstlisting}
; Call function at address in a
LOAD ah, #0x0040
LOAD al, #0x0000              ; a = 0x00400000
LJSR a                        ; Call 0x00400000

; Call with offset
LJSR a, #function_offset      ; p = a + function_offset

; Call with register displacement
LOAD r1, #8
LJSR a, r1                    ; p = a + r1
\end{lstlisting}

\end{instructionbox}

\section{Control Flow Patterns}

\subsection{If-Then-Else}
\begin{lstlisting}
; if (r1 > r2) { ... } else { ... }
CMPR r1, r2
BRAN|<= else_part
; Then part
; ...
BRAN end
else_part:
; Else part
; ...
end:
\end{lstlisting}

\subsection{While Loop}
\begin{lstlisting}
; while (r1 > 0) { ... }
while_start:
    CMPI r1, #0
    BRAN|<= while_end
    ; Loop body
    ; ...
    BRAN while_start
while_end:
\end{lstlisting}

\subsection{Switch/Jump Table}
\begin{lstlisting}
; Switch on r1 (0-3)
LOAD r2, r1                   ; r2 = r1
ADDI r2, #0, LSL #1           ; r2 = (r2 << 1) = r2 * 2 (word size)
LDEA p, (r2, p)               ; Jump via table

jump_table:
    DW case0 - jump_table
    DW case1 - jump_table
    DW case2 - jump_table
    DW case3 - jump_table
\end{lstlisting}
